\section{A verified composition of the two 239-state reductions}

The starting point for this section is the boundary-reassociated 239-state
machine, denoted $B$, with inherited start state $95$, dispatcher $96$, and
separate halt state $254$.  Its table omits state $218$ and retains state
$169$.  The coordinated-helper reduction from the other 239-state
construction acts on a different set of states.  Applying that reduction
to $B$ produces a machine $C$ with 238 working states.  Compatibility is a
mathematical assertion that needs verification: changes to the decoder can
change the contexts in which a shared state is visited.  The checks below
establish the required compatibility for this exact pair of machines.

Define the control map
\[
 \phi(q)=\begin{cases}80,&q=169,\\q,&q\ne169.\end{cases}
\]
Delete the two rows of state $169$, apply $\phi$ to every target, and replace
the two retained rows
\[
 C(80,0)=(0,L,106),\qquad C(91,1)=(1,L,167).
\]
The only redirected incoming edge is $147/0$, whose target changes from
$169$ to $80$.  Thus the construction changes three retained rows, removes
two rows, and leaves the other 473 retained rows literally unchanged.
The new machine still has alphabet $\{0,1\}$ and a two-way infinite tape;
every transition moves exactly one cell left or right.  Its entire blank
initializer is included in the 238-state count.

\subsection{The predecessor invariant}

Write $t(i)$ for the tape symbol at position $i$, and let $h$ be the head
position.  The following implication is an inductive invariant of $B$ from
state $95$, even if its initial tape is arbitrary:
\[
\begin{array}{c|l}
\text{state}&\text{forced symbols}\\\hline
67&t(h-1)=0\\
80&t(h-2)=0,\quad t(h-1)=1\\
91&t(h-3)=0,\quad t(h-2)=t(h-1)=1\\
169&t(h-1)=0.
\end{array}
\]
Here and below, an implication associated with one of these states makes
no demand in any other state.

The proof uses the complete incoming-edge inventory.  State $67$ can be
entered only from $52/0$ or $124/0$, both writing zero and moving right.
State $80$ can be entered only from $67/1$, writing one and moving right;
the preceding zero from the invariant of state $67$ then lies two cells
left.  State $91$ can be entered only from $80/1$, writing one and moving
right, which shifts the forced pattern to $011$.  State $169$ can be
entered only from $147/0$, writing zero and moving right.  None of these
conditions is imposed at the start state $95$.  These observations prove
initiality and preservation by every transition.

The finite-window replay independently checks this invariant on all
121,024 admitted combinations of a working state and a nine-cell Boolean
window.  The predecessor proof explains why those calculations apply to
an arbitrary exterior tape; the machine reads and writes only the scanned
cell, and preservation inspects only the explicitly listed neighbors.

\subsection{Four positive blocks on an open tape}

The additional premise required by the reduction is
\begin{equation}\label{eq:boundary169-right-guard}
 q=169\ \wedge\ t(h)=1\quad\Longrightarrow\quad t(h+1)=1.
\end{equation}
This premise does not follow from the preceding incoming-edge argument.
It is proved for the blank computation by the unbounded instruction-frame
argument in the next subsection.

With that premise, each exceptional step of $B$ has the following exact
simulation in $C$:
\[
\begin{array}{c|c|l|r}
\text{old state}&\text{read}&\text{new state path}&\text{steps}\\\hline
80&0&80,106,92,107,125,25&5\\
169&0&80,106,93,57&3\\
169&1&80,91,167,151,173,56&5\\
91&1&91,167,151,24,36,54,71,82,94,50&9.
\end{array}
\]
The final entry in each path is the target configuration's state; it does
not contribute another transition.  Each displayed block agrees with the
old single step in the entire final tape, head position, and state after
applying $\phi$.

For $80/0$, the forced symbol immediately left is one.  The new head
positions, relative to its starting point, are
$-1,0,-1,0,1$; all inspected symbols are the forced left one and the
scanned zero.  Every write preserves its cell, and the final state is $25$.
For $169/0$, the forced left zero selects $106/0$ instead.  The head follows
$-1,0,1$, all writes preserve zeros, and the final state is $57$.

For $169/1$, the new first step temporarily leaves the scanned one in place
and moves right to state $91$.  Premise~\eqref{eq:boundary169-right-guard}
forces this right neighbor to be one.  The head positions then become
$1,0,-1,0,1$.  The forced left zero selects $151/0$, and $173/1$ erases
the original scanned one in the final step.  The right neighbor is
preserved.  The final effect is precisely the old $169/1$ transition,
which writes zero, moves right, and enters $56$.

For $91/1$, the pattern at positions $h-3,h-2,h-1,h$ is $0111$.
The head positions are
$-1,-2,-3,-2,-1,0,-1,0,-1$.
Every transition preserves its symbol until $94/1$ erases the original
scanned one and moves left to $50$.  This is the complete effect of the
old $91/1$ transition.  No part of this argument inspects a cell outside
the forced pattern.

The independent open-tape checker represents only the forced cells and
rejects any trace attempting to read another cell.  It therefore proves
these identities with no assignments to exterior cells.  A separate
finite-footprint checker obtains the same result on all 22 admitted
assignments to the window from $h-3$ through $h+1$.  The remaining 474 old
transition rows commute literally with $\phi$ in one step.  No displayed
block visits the halt state in its interior.

\subsection{The unbounded right-neighbor guard}

The decoder source is the exact boundary-reassociated list of 6,455
instruction sites.  The guard audit keeps this source and its 239-state
table fixed and rederives the four arbitrary-list traversal lemmas from
the literal table.  It then checks every instruction site, splitting
conditional decrements into their zero and positive cases.  This gives
8,386 symbolic cases.  The completed fresh audit uses 2,247,513 symbolic
rules and covers all 478 rows of the old table.  It explicitly checks 32
state-$169$/read-one visits; the smallest constant term of their one-run
lengths is 23.  The resulting bound is stronger than the two consecutive
ones needed by the local reduction.  The minimum certified constant part
of an instruction duration is 13.

Each case has nine independent natural register values, an arbitrary
additional gap, an arbitrary first spare value, and an arbitrary finite
remaining spare list.  Unary run lengths are affine forms in independent
natural variables.  At every literal visit to old state $169$ scanning a
one, the audit normalizes the run representation and checks that the
constant term of the current one-run length is at least two.  Since all
variable coefficients are nonnegative, the right neighbor is then one
for every assignment of those parameters.

The audit also checks that state $169$ is absent from the interior traces
of all accelerated list-traversal lemmas, is not the entry state of any
such lemma, and has no self-loop on either symbol.  Consequently neither
list induction nor unary-loop induction can conceal a state-$169$ visit.
Every applicable visit is explicitly inspected.  Full instruction
endpoints and positive duration checks are retained unchanged from the
underlying mathematical runtime checker.  The finite cap used to build
a symbolic derivation bounds the proof construction, not the values of
the registers, the length of the spare list, or the actual computation.

The exact blank prefix supplies the remaining initialization case.
All 409,692 old checkpoints, including both endpoints of the 409,691-step
initialization, agree with the corresponding complete new configurations
under $\phi$.  The six old initialization visits to $169/1$ satisfy
\eqref{eq:boundary169-right-guard}.  The counts of the four exceptional
rows $(80/0,169/0,169/1,91/1)$ are $(18,12,6,3)$.  Therefore the exact new
initialization time is
\[
409691+4\cdot18+2\cdot12+4\cdot6+8\cdot3=409835.
\]
The final head is at $-5$, the dispatcher is $96$, and the entire tape
has the same 351 one-cells as the old complete startup frame; every other
cell is zero.  Concatenating the unbounded instruction cases from that
frame establishes the right-neighbor guard throughout the old blank
computation.

\subsection{Why positive simulation gives halting equivalence}

For completeness, the simulation argument can be stated without a
bounded-execution assumption.  Let $c_0,c_1,\ldots$ be the old blank
computation, up to a halt if there is one.  Keep tape and head unchanged
in the configuration map $\Phi$ and replace its control state by $\phi$.
Each old transition has a certified block length $\ell(c_i)$ between
one and nine, with
\[
 C^{\ell(c_i)}(\Phi(c_i))=\Phi(c_{i+1}).
\]
Set $k_0=0$ and $k_{i+1}=k_i+\ell(c_i)$.  Induction gives
$C^{k_i}(\Phi(c_0))=\Phi(c_i)$, and no corresponding intermediate
configuration halts.  If the old machine halts, the new machine reaches
its mapped halt.  If the old machine has infinitely many steps, then
$k_i\ge i$ is unbounded; the positive blocks cover the entire new
computation, which therefore never halts.  Determinism yields both
directions of halting equivalence.

All arguments here concern natural numbers, finite strings, finite
lists, and their inductions, and are formalizable in ZFC.  Composing this
equivalence with the inherited calibration of the boundary source gives
the requested equivalence between nonhalting of the exact 238-state
machine and $\operatorname{Con}(\mathrm{ZFC})$.  This is an ordinary
mathematical proof with executable finite and unbounded symbolic
certificates.  It does not constitute a newly kernel-checked theorem for
the 238-state table, nor does it assert that ZFC proves its own
consistency or that 238 is a minimal state count.
